% !TeX root = ../main.tex

\documentclass[../main.tex]{subfiles}
\begin{document}

\begin{center}
    \Large{\textbf{TÓM TẮT NỘI DUNG}}\\
\end{center}
\vspace{1cm}

Đồ án nghiên cứu và đề xuất phương án công nghệ xử lý chất thải rắn sinh hoạt (CTRSH) công suất 120 tấn/ngày đêm cho Nhà máy xử lý CTRSH thành phố Hội An. Xuất phát từ đặc tính rác đầu vào đặc thù của địa phương du lịch, phương án nhiệt phân lò quay liên tục kết hợp chưng cất dầu và phát điện bằng động cơ diesel được lựa chọn thay thế cho dây chuyền khí hóa -- tuabin khí ban đầu.

Nội dung chính của đồ án bao gồm: luận chứng lựa chọn công nghệ; tính toán cân bằng vật chất và cân bằng năng lượng cho toàn hệ thống; thiết kế sơ bộ các thiết bị chính; phân tích hiệu quả kinh tế thông qua CAPEX, OPEX, NPV và IRR; đồng thời đánh giá các yêu cầu kiểm soát môi trường và bảo đảm an toàn vận hành.

Kết quả tính toán xác nhận phần cháy được của CTRSH sau tiền xử lý tạo ra nhiên liệu từ rác (RDF) có nhiệt trị ổn định, từ đó hình thành các dòng sản phẩm gồm dầu nhiệt phân, khí không ngưng và than carbon phục vụ thu hồi năng lượng. Sản lượng điện phát ra đủ để đáp ứng toàn bộ nhu cầu vận hành nhà máy và có phần dư xuất lưới. Phương án cho thấy tỷ lệ chôn lấp cuối cùng giảm đáng kể so với mô hình xử lý hiện hành.

Kết quả của đồ án là một phương án thiết kế sơ bộ có tính nhất quán về kỹ thuật và phù hợp với điều kiện thực tế của dự án, làm cơ sở tham khảo cho thiết kế chi tiết và triển khai các mô hình xử lý CTRSH tương tự tại Việt Nam.

% \begin{flushright}
% Sinh viên thực hiện\\
% \begin{tabular}{@{}c@{}}
% \textit{(Ký và ghi rõ họ tên)}
% \end{tabular}
% \end{flushright}

\end{document}
